\section{Examples and application adapters}
\label{app:additional-examples}

The worked grid in Section \ref{ch:worked} is the shortest route through the
entire calculation, but it is not the package's only example. This appendix
documents the repository's ten examples, including the
external-data builders whose raw inputs remain outside Git. The examples serve
different purposes. The toy, Braess, and multimodal urban networks isolate
specific parts of the calculation. Sioux Falls adapts a familiar traffic
network, while Seattle and Westeros apply the two spatial models to larger
geographies. A numerical result appears only when the corresponding input
passes the package's accounting and numerical checks.

\subsection{Coverage and status}

\begin{longtable}{L{0.15\textwidth}L{0.18\textwidth}L{0.29\textwidth}L{0.28\textwidth}}
\toprule
Example & Spatial model & Construction & Status and role \\
\midrule
\endhead
Three-node toy & Economic geography & Hand-written labor, income, road, rail,
and terminal records & Self-contained API and sensitivity smoke test \\
Multimodal grid & Economic geography & Programmatically generated \(5\times5\)
road grid with a central transit spine & Self-contained worked example in
Section \ref{ch:worked} \\
Braess-style network & Economic geography & Hand-written four-node, two-route
network with a connector & Self-contained route-closure test \\
Cow network & Economic geography & Original 30-node three-dimensional
geography; optional external PLY mesh & Self-contained ranking example; large
mesh is on demand \\
Sioux Falls & Economic geography & Pinned TNTP network, trips, assigned flows,
and local BPR elasticities & Verified external-data benchmark adaptation \\
Three-location urban toy & Urban commuting & Hand-written residence,
workplace, OD, and road-flow records & Self-contained one-mode regression
oracle \\
Multimodal urban toy & Urban commuting & Four locations, road and transit,
alternative routes, and terminal congestion & Self-contained shared-transport
test \\
Seattle road diagnostic & Urban commuting & Published LODES and HPMS/HERE
inputs from the Allen--Arkolakis replication & Rejected by the accounting
gate; no welfare result is reported \\
Seattle multimodal & Urban commuting & LODES OD, ACS mode shares, and
hash-pinned 2017 GTFS service & Verified candidate application, not an
externally calibrated appraisal \\
Westeros & Economic geography & Public ArcGIS geography with synthetic
activity, trade, and traffic & Verified illustrative calculation; not an
empirical claim \\
\bottomrule
\end{longtable}

The maps, scatters, and ranked tables below follow one design rule. Map color
records the Extended-approach welfare elasticity, line width rises with its
absolute value, and node area records economic activity. Each scatter uses the
same scale on both axes, so the dashed line marks equality between the
Traditional and Extended approaches. The figures generated from external data
are frozen to input and output hashes. Rebuilding them requires the source data,
but compiling this guide does not.

\subsection{Self-contained economic-geography examples}

These examples keep their model inputs in the repository and differ in network
structure and the economic response they are designed to expose. The smallest
case checks the API; the remaining cases add route, modal, and congestion
comparisons.

\subsubsection{Three-node API toy}

The smallest generic example has three locations and three reciprocal physical
links. Road is available on all three links, while rail is available on two links
and loads an endpoint terminal at each end. The input files record positive
labor and income at each location and model-ready value flows by edge and
mode. The configuration combines edge-local road congestion with rail-terminal
congestion and evaluates a one-percent bidirectional road improvement. Because
the files are short enough to inspect directly, this example is the package's
first check of parsing, normalization, modal aggregation, terminal identifiers,
policy-unit handling, decomposition, and the seven sensitivity paths.

\begin{lstlisting}[language=bash]
julia --project=. bin/tnw.jl validate examples/toy/config.toml
julia --project=. bin/tnw.jl decompose examples/toy/config.toml
\end{lstlisting}

The example is original to this package. Its coordinates are only plotting
metadata, and its flows are economic value shares rather than measurements
from the three cities named by those coordinates.

\subsubsection{Multimodal grid}

The grid used throughout the guide is generated by
\code{examples/grid\_multimodal/build\_inputs.jl}. Reciprocal road links join
adjacent locations on a \(5\times5\) lattice. Transit is available on the four
links along the middle row, and explicit terminal identifiers load the
terminal-congestion extension. Symmetric directional flows and positive
location activity are chosen together so that the recursive-flow identities
hold without balancing inside the estimator. Section \ref{ch:worked} reports
the network map, Traditional--Extended comparison, ranked links, and exact
decomposition. The example and all of its displayed assets are original and
fully distributable.

\subsubsection{Braess-style routing network}

The Braess-style example uses the four-node topology associated with the
classic traffic-planning paradox \citep{Braess1968}. Two paths connect the
outer locations, and a cross-link joins the intermediate nodes. The package
does not solve a Wardrop assignment, so the example does not claim to reproduce
Braess's paradox. Instead, reciprocal model-ready value flows are placed on
the five physical links. The full flexible-route response can then be compared
with \code{FR}, which holds baseline physical-edge use fixed. This construction
isolates the route term while retaining the economic-geography equilibrium.

\begin{figure}[H]
\centering
\includegraphics[width=0.98\textwidth]{braess-comparison.pdf}
\caption{Braess-style routing example. The left panel maps the
Extended-approach elasticity. The right panel compares the Traditional and
Extended approaches on common axes.}
\label{fig:braess-comparison}
\end{figure}

\begin{table}[H]
\centering
\small
\caption{Physical links in the Braess-style example}
\label{tab:braess-top-links}
% Generated by plots/example_assets.py; do not edit.
\begin{tabular}{lrrrr}
\toprule
Link & Traditional & Extended & Traditional rank & Extended rank \\
\midrule
AB & 0.160000 & 0.103643 & 1 & 1 \\
BD & 0.160000 & 0.103643 & 2 & 2 \\
AC & 0.100000 & 0.067615 & 3 & 3 \\
CD & 0.100000 & 0.067615 & 4 & 4 \\
BC & 0.040000 & 0.025182 & 5 & 5 \\
\bottomrule
\end{tabular}

\end{table}

The route comparison is nonzero because the connector and the two outer paths
give traffic more than one way to reach the destination. All five
Extended-approach points lie below the equality line in this calibration. That
ordering is a result of the chosen externalities, congestion, and equilibrium
response; it is not imposed by the routing operator.

\clearpage
\subsubsection{Cow network and full mesh}

The 30-node cow is an original synthetic geography. Nodes and links form a
three-dimensional body, head, legs, and two alternative neck routes. Positive
activity is concentrated in the body and head. Model-ready balanced flows then
give the neck a clear economic role: a link can matter because it carries
traffic, because it changes access between the two high-activity regions, or
because it changes how traffic uses the alternative route. The elevation
coordinate affects only the visualization.

\begin{figure}[H]
\centering
\includegraphics[width=0.98\textwidth]{cow-comparison.pdf}
\caption{Cow-shaped economic geography. Both panels use the common example
design; the three-dimensional rendering supplied by the package uses the same
welfare scale.}
\label{fig:cow-comparison}
\end{figure}

\begin{table}[H]
\centering
\small
\caption{Highest-ranked physical links in the cow-shaped network}
\label{tab:cow-top-links}
% Generated by plots/example_assets.py; do not edit.
\begin{tabular}{lrrrr}
\toprule
Link & Traditional & Extended & Traditional rank & Extended rank \\
\midrule
C1-C2 & 0.120000 & 0.049769 & 1 & 1 \\
H3-H4 & 0.050000 & 0.034154 & 21 & 2 \\
H2-H3 & 0.050000 & 0.034081 & 20 & 3 \\
B2-C1 & 0.090000 & 0.034050 & 2 & 4 \\
B4-C2 & 0.090000 & 0.033947 & 3 & 5 \\
H1-N1 & 0.090000 & 0.031296 & 5 & 6 \\
B8-C1 & 0.080000 & 0.029924 & 8 & 7 \\
B3-B4 & 0.080000 & 0.029821 & 6 & 8 \\
\bottomrule
\end{tabular}

\end{table}

Only \(\GuideCowTopOverlap\) of the Traditional top
\(\GuideCowTopCount\) links remain in the Extended top
\(\GuideCowTopCount\), and the largest rank movement is
\(\GuideCowMaximumRankShift\) positions. The closure results contain nonzero
congestion and route terms, which are two channels through which the Extended
ranking can depart from the traffic ranking. The calculation does not assign
each individual rank change to a single channel.

The optional high-resolution version begins with John Burkardt's 2,903-vertex
PLY cow \citep{BurkardtCowPLY}. The builder assigns positive labor and income
to all 2,903 mesh vertices, turns each mesh edge into a reciprocal road link, and
uses symmetric matrix scaling to construct balanced value flows. The resulting
network has 8,706 physical links and 17,412 directed arcs. Its sparse
edge-local calculation evaluates the same primitive-cost derivative, but it
does not build the dense fixed-route and fixed-mode incidence objects required
for \code{FR} and \code{FM}. The PLY file is downloaded and hash-checked on
demand under its stated GNU LGPL license; it is not stored in the repository.

\subsection{Urban commuting examples}

The urban examples retain the common transport equations but replace labor and
income shares with residence and workplace margins. The one-mode oracle first
checks the urban welfare derivative in isolation; the multimodal network then
adds the shared route and mode machinery.

\subsubsection{Three-location regression oracle}

The one-mode urban toy follows the residence--workplace structure of
\citet{AA_2022restud}. Its node file records distinct residence and workplace
masses, and its edge file records directed commuter traffic. The data are
constructed to satisfy the urban flow identity exactly. With spatial
externalities and congestion set to zero, the urban derivative collapses to
commuter traffic. The test suite also solves the nonlinear exact-hat system
around this baseline and compares its central finite difference with the
analytical derivative. The example therefore checks the urban Jacobian and
welfare projection independently of modal or terminal terms.

\subsubsection{Multimodal urban network}

The four-location urban example retains the same residence--workplace
equilibrium but adds the transport block used by the economic-geography model.
Road and transit operate on all ten directed edges, two routes connect the main
origin and destination, road traffic creates edge congestion, and transit
loads endpoint terminals. Residence shares, workplace shares, OD flows, and
edge-mode flows are generated from one contractive route kernel, so no
balancing occurs during estimation.

\begin{figure}[H]
\centering
\includegraphics[width=0.98\textwidth]{urban-comparison.pdf}
\caption{Multimodal urban example. Dashed blue segments identify links with
transit service, and blue rings identify terminal locations.}
\label{fig:urban-comparison}
\end{figure}

\begin{table}[H]
\centering
\small
\caption{Physical links in the multimodal urban example}
\label{tab:urban-top-links}
% Generated by plots/example_assets.py; do not edit.
\begin{tabular}{lrrrr}
\toprule
Link & Traditional & Extended & Traditional rank & Extended rank \\
\midrule
BD & 0.052417 & 0.035007 & 1 & 1 \\
AB & 0.047878 & 0.032458 & 2 & 2 \\
AC & 0.035802 & 0.023888 & 3 & 3 \\
CD & 0.024323 & 0.015681 & 4 & 4 \\
BC & 0.006896 & 0.004512 & 5 & 5 \\
\bottomrule
\end{tabular}

\end{table}

The example changes the spatial residuals and welfare row but reuses the
economic-geography route, mode, congestion, and pass-through operators. Its
no-congestion derivative equals the urban traffic statistic to within
\(\GuideUrbanHultenCollapseTolerance\). This is the package's direct numerical
check that multimodality belongs to the shared transport block rather than to
a separate urban transport model.

\subsection{Sioux Falls benchmark adaptation}

Sioux Falls is a 24-node test network introduced in early traffic-assignment
work and widely used as a small computational benchmark
\citep{LeBlancMorlokPierskalla1975}. The package pins the network, assigned
flow, node-coordinate, and trip files in the Transportation Networks for
Research repository \citep{TransportationNetworksSiouxFalls}. Their stated
terms limit the files to academic research, so the builder downloads them on
demand and verifies the pinned commit and all four hashes.

The benchmark files are not immediately model-ready. OD productions and
attractions differ slightly, and the assigned directional flows reproduce
those differences. The builder averages production with attraction and each
pair of reciprocal assigned flows. The balanced OD margin becomes both labor
and income, while balanced link flow is divided by total OD trips. It records
the original imbalance and verifies exact conservation after this
transformation. The extended configuration evaluates each BPR cost function
at the assigned baseline and uses its local log derivative,
\[
 \lambda_e=
 \frac{p_e B_e(v_e/c_e)^{p_e}}
      {1+B_e(v_e/c_e)^{p_e}},
\]
as the link-specific congestion elasticity. This is an adaptation of the
assigned baseline to recursive routing, not a new solution of the original
Wardrop problem.

\begin{figure}[H]
\centering
\includegraphics[width=0.98\textwidth]{sioux-comparison.pdf}
\caption{Sioux Falls benchmark adaptation. The calculation uses
\GuideSiouxNodeCount\ nodes, \GuideSiouxDirectedArcCount\ directed arcs, and
\GuideSiouxPhysicalLinkCount\ reciprocal physical links.}
\label{fig:sioux-comparison}
\end{figure}

\begin{table}[H]
\centering
\small
\caption{Highest-ranked physical links in the Sioux Falls adaptation}
\label{tab:sioux-top-links}
% Generated by plots/example_assets.py; do not edit.
\begin{tabular}{lrrrr}
\toprule
Link & Traditional & Extended & Traditional rank & Extended rank \\
\midrule
12-13 & 0.068403 & 0.042072 & 13 & 1 \\
1-3 & 0.044963 & 0.037363 & 30 & 2 \\
3-12 & 0.055452 & 0.031508 & 19 & 3 \\
1-2 & 0.024996 & 0.023084 & 38 & 4 \\
7-18 & 0.087767 & 0.019054 & 8 & 5 \\
3-4 & 0.077751 & 0.016411 & 11 & 6 \\
18-20 & 0.105295 & 0.016382 & 4 & 7 \\
16-18 & 0.084891 & 0.015318 & 10 & 8 \\
\bottomrule
\end{tabular}

\end{table}

The Traditional and Extended rankings differ substantially: their Pearson
correlation is \(\GuideSiouxPearson\), and their rank correlation is
\(\GuideSiouxSpearman\). The mean elasticity falls from
\(\GuideSiouxMeanTraditional\) to \(\GuideSiouxMeanExtended\). Link-specific
BPR derivatives enter the Extended measure but not the Traditional measure,
so their variation is one reason the rankings diverge. The calculation is not
a validated welfare estimate for Sioux Falls. It shows how the package
transforms an assigned-flow traffic benchmark before applying the welfare
formula.

\subsection{Seattle road diagnostic and multimodal candidate}

Seattle provides the larger urban application, but its published road inputs
and the candidate multimodal construction have different evidentiary status.
The distinction matters before interpreting either set of welfare results.

\subsubsection{The published road inputs remain a diagnostic}

The first Seattle adapter preserves the 217-location inputs used by
\citet{AA_2022restud}: 2017 LODES residence--workplace flows and 2016 road
traffic and travel-time inputs. It verifies the archive hashes, 217 nodes,
1,384 directed road arcs, and 692 reciprocal links. It then fails validation.
LODES counts commuters, whereas HPMS annual average daily traffic counts road
vehicles from a different population. The two sources imply a maximum
normalized recursive-stock disagreement of about \(0.108\). The package does
not silently balance that difference, and the guide therefore reports no
welfare map, scatter, or ranking from this rejected baseline.

\subsubsection{Construction of the multimodal baseline}

The multimodal candidate instead constructs all transport flows from one
commuter population. The complete 2017 LODES matrix supplies OD commuter
masses and its margins supply residence and workplace activity
\citep{CensusLODES}. ACS table B08301 supplies each residential origin's
public-transit share \citep{CensusACS2017B08301}. A hash-pinned June 2017 King
County GTFS feed supplies service, paths, and bus, rail/streetcar, and ferry
labels \citep{KingCountyGTFS2017}. Published Seattle road travel times supply
road paths. The builder routes each OD-mode flow over these networks, then
checks that the resulting edge-mode traffic reproduces the OD margins.

This construction does not turn schedules into observed ridership. GTFS
identifies feasible services and scheduled travel, while ACS identifies mode
shares at residential origins. The 2017 Metro System Evaluation provides a
separate route-activity comparison and does not determine the OD-to-route
allocation \citep{KingCountyMetro2017Evaluation}. The modal elasticity
\(\eta=1.099\) is transferred from the economic-geography application and is
reported with sensitivity rather than described as a Seattle estimate.

\begin{figure}[p]
\centering
\includegraphics[width=0.84\textwidth]{seattle-network.pdf}
\caption{Observed 2017 Seattle networks and the constructed multimodal model
network. Road, bus, rail/streetcar, and ferry use distinct colors. Node area
records residence plus workplace commuter mass.}
\label{fig:seattle-network}
\end{figure}

\begin{figure}[p]
\centering
\includegraphics[height=0.84\textheight]{seattle-mode-map.pdf}
\caption{Seattle road and transit policy support. Color identifies the mode,
and line width records the Extended-approach welfare gain from a one-percent
improvement. Dashed segments have a negative Extended-approach effect.}
\label{fig:seattle-mode-map}
\end{figure}

\begin{figure}[p]
\centering
\includegraphics[width=0.92\textwidth]{seattle-welfare-map.pdf}
\caption{Mode-specific Seattle welfare maps on a common scale. Each panel
reports the Extended-approach gain from a one-percent edge-mode improvement,
not baseline traffic or ridership.}
\label{fig:seattle-welfare-map}
\end{figure}

\begin{figure}[p]
\centering
\includegraphics[width=0.88\textwidth]{seattle-scatter.pdf}
\caption{Traditional and Extended approaches for Seattle transit
improvements. The four panels share both axes and include the equality line.}
\label{fig:seattle-scatter}
\end{figure}

\begin{table}[p]
\centering
\scriptsize
\caption{Selected high-ranked Seattle corridors by mode}
\label{tab:seattle-top-links}
% Generated by scripts/build_practitioner_external_example_assets.py.
\begin{tabular}{llrrrr}
\toprule
Corridor & Mode & Traditional & Extended & Traditional rank & Extended rank \\
\midrule
114\_121 & road & 0.139194 & 0.064713 & 2 & 1 \\
100\_95 & road & 0.139007 & 0.052276 & 3 & 2 \\
107\_114 & road & 0.124682 & 0.050706 & 4 & 3 \\
100\_95 & bus & 0.061214 & 0.024548 & 1 & 1 \\
100\_106 & bus & 0.037927 & 0.021926 & 2 & 2 \\
71\_95 & bus & 0.017758 & 0.020081 & 10 & 3 \\
102\_123 & rail & 0.020337 & 0.012086 & 1 & 1 \\
49\_67 & rail & 0.010131 & 0.007669 & 2 & 2 \\
28\_49 & rail & 0.004571 & 0.002744 & 3 & 3 \\
89\_95 & ferry & 0.000888 & 0.001377 & 1 & 1 \\
\bottomrule
\end{tabular}

\end{table}

The candidate evaluates \GuideSeattleRoadArcCount\ directed road,
\GuideSeattleBusArcCount\ bus, \GuideSeattleRailArcCount\ rail/streetcar, and
\GuideSeattleFerryArcCount\ ferry interventions on the same
\GuideSeattleNodeCount-location equilibrium. Mean Extended-approach gains are
\(\GuideSeattleBusMeanExtended\) percent for bus,
\(\GuideSeattleRailMeanExtended\) percent for rail/streetcar, and
\(\GuideSeattleFerryMeanExtended\) percent for ferry. The bus and rail rank
correlations with the Traditional approach are \(\GuideSeattleBusSpearman\)
and \(\GuideSeattleRailSpearman\), respectively. The maximum error in the
stratified nonlinear finite-difference checks is
\(\GuideSeattleFDMaximum\). The accounting identities, deterministic rebuild,
and finite-difference check verify the calculation for the stated inputs and
calibration. They do not supply route-level ridership, a Seattle estimate
of modal substitution, or crowding parameters. The application is therefore
an illustration rather than a complete transit appraisal.

\subsection{Westeros economic geography}

The Westeros example uses the public locations, roads, and continent polygon
in CreativeCarto's ArcGIS item \citep{CreativeCartoWesteros}. Those layers
contain geography but no economic activity or traffic. The current example
therefore uses the economic-geography model, not the urban commuting model.
The builder hash-checks the source responses, connects road components within a
declared 10-kilometer threshold, retains settlements within 50 kilometers of
the connected road system, and forms a sparse settlement graph from road
distances.

The economic baseline is synthetic and fully disclosed. Location types receive
positive common labor and income weights. A symmetric, doubly constrained
gravity calculation assigns bilateral trade, and shortest paths route that
trade over the reduced road graph. The builder verifies the trade margins and
edge-flow accounting before writing the package inputs. The hash-pinned build has
\GuideWesterosNodeCount\ locations, \GuideWesterosDirectedArcCount\ directed
arcs, and \GuideWesterosPhysicalLinkCount\ physical links.

\clearpage
\begin{figure}[H]
\centering
\includegraphics[height=0.82\textheight]{westeros-world-map.pdf}
\caption{Synthetic Westeros economic geography. The land outline and faint
road network come from the CreativeCarto ArcGIS layers
\citep{CreativeCartoWesteros}. Colored lines trace the evaluated model links
along those roads; their color and width record the Extended-approach welfare
elasticity. Node area records synthetic economic activity.}
\label{fig:westeros-world-map}
\end{figure}

\clearpage
\begin{figure}[H]
\centering
\includegraphics[width=0.78\textwidth]{westeros-scatter.pdf}
\caption{Traditional and Extended welfare elasticities in the synthetic
Westeros calculation. Both axes use the same scale, and selected links are
identified by their endpoint names.}
\label{fig:westeros-scatter}
\end{figure}

\begin{table}[H]
\centering
\scriptsize
\caption{Highest-ranked physical links in the Westeros example}
\label{tab:westeros-top-links}
% Generated by plots/westeros_example.py; do not edit.
\begin{tabular}{L{0.31\textwidth}rrrr}
\toprule
Link & Traditional & Extended & Trad. rank & Ext. rank \\
\midrule
Ashemark--Casterly Rock & 0.083333 & 0.014754 & 9 & 1 \\
Felwood--King's Landing & 0.094063 & 0.013233 & 6 & 2 \\
Riverrun--Stone Hedge & 0.092104 & 0.012854 & 7 & 3 \\
Hayford Castle--King's Landing & 0.103907 & 0.012797 & 4 & 4 \\
Goldengrove--Highgarden & 0.062883 & 0.012383 & 13 & 5 \\
Riverrun--Wayfarer's Rest & 0.083861 & 0.012160 & 8 & 6 \\
Casterly Rock--Lannisport & 0.047727 & 0.012054 & 20 & 7 \\
Darry--Maidenpool & 0.147725 & 0.010379 & 1 & 8 \\
\bottomrule
\end{tabular}

\end{table}

The two welfare measures have a Pearson correlation of
\(\GuideWesterosPearson\) and a rank correlation of
\(\GuideWesterosSpearman\). Their means are
\(\GuideWesterosMeanTraditional\) and
\(\GuideWesterosMeanExtended\), respectively. These numbers describe how the
package behaves on an irregular geography; they do not measure a fictional
economy. The ArcGIS item did not state a reuse license when the source manifest
was recorded. Raw responses and generated model inputs consequently remain
outside Git, and any redistributed display should retain the source
attribution.

\subsection{Comparison of mechanisms across the self-contained examples}

Table \ref{tab:additional-example-mechanisms} reports mean differences across
the exact closure cases for the four self-contained examples used in the
common asset build. ``Congestion'' is
\(\overline E_F^r-\overline E_{NC}^r\), ``Mode flexibility'' is
\(\overline E_F^r-\overline E_{FM}^r\), and ``Route flexibility'' is
\(\overline E_F^r-\overline E_{FR}^r\). The entries are signed elasticity
differences rather than percentage shares, so they may have opposite signs.

\begin{table}[H]
\centering
\small
\caption{Mean welfare differences across self-contained examples}
\label{tab:additional-example-mechanisms}
% Generated by scripts/build_practitioner_guide_assets.jl; do not edit.
\begin{tabular}{lrrrr}
\toprule
Example & Congestion & Mode flexibility & Route flexibility & Top-$k$ overlap \\
\midrule
Multimodal grid (EG) & 0.001685 & -0.000001 & -0.000341 & 5/5 \\
Braess-style routes (EG) & 0.012657 & 0.000000 & -0.002924 & 1/1 \\
Cow network (EG) & 0.015092 & 0.000000 & -0.007406 & 3/5 \\
Multimodal urban (U) & 0.002522 & -0.000003 & -0.000359 & 1/1 \\
\bottomrule
\end{tabular}

\end{table}

Congestion raises the realized-cost derivative in all four included
calibrations, while route flexibility lowers it relative to fixed baseline
routes. Modal flexibility is small in these particular synthetic baselines.
That pattern is an output, not a maintained restriction. The examples are
most useful when read together: the grid combines all channels, the
Braess-style topology isolates alternative routes, the cow creates uneven
market access on an irregular network, and the urban example changes the
spatial equilibrium while retaining the same transport operators.
